\section{Fixed-target extremizers and the conserved resource}
\label{v9:core:section}

We optimize the excess above a fixed candidate constant. This gives a
variational problem whose coefficient is known before any acquisition
step. The positive excess is retained through the contact calculation
and the conserved identity; it becomes a nonnegative term in the final
resource. No upper bound for the maximizing normalized score is needed.
Throughout, sums have
independent identically distributed summands with mean zero, variance one
and finite third absolute moment $\rho$. Fix $n\ge2$, $c>0$, and a real
raw sum threshold $H$. Set
\begin{gather*}
 z=H/\sqrt n,\quad F=\Pp(S_n\le H),\quad
 a=\sqrt n\{F-\Phi(z)\}-c\rho,\\
 \varphi=\phi(z),\quad\beta=\E[X|X|],\quad\gamma=\E X^3.
\end{gather*}
In particular $a$ is not divided by $\rho$.

\subsection{Attainment, contamination and the entire closed support}

\begin{proposition}\label{v9:core:attainment}
For each fixed $n$ and $c>0$, the supremum of $a$ over these laws and
thresholds is attained, and the closed-CDF discrepancy at every attaining
pair is positive.
At every maximizing pair,
\begin{align}
 I(y)={}&n^{3/2}\{h(y)-F\}+ny\varphi
       +\tfrac12H\varphi(y^2-1)\notag\\
 &+\tfrac32c\rho(y^2-1)-c\{|y|^3-\rho-3y\beta\}\le0,
 \qquad h(y)=\Pp(S_{n-1}\le H-y),\label{v9:core:contact}
\end{align}
with equality on the whole closed support. That support is compact.
Neither statement assumes $a\ge0$ or $H\ge0$.
\end{proposition}
\begin{proof}
Write $A$ for the supremum. Since $\rho\ge1$ and the discrepancy is
at most one, $A\le\sqrt n-c$. Symmetric Bernoulli evaluated at $H=n$
gives $A\ge\sqrt n(1-\Phi(\sqrt n))-c>-c$.
For a maximizing sequence, $a_j\to A$ and
$c\rho_j\le\sqrt n-a_j$ give a common third-moment ceiling.
The discrepancies satisfy $\sqrt n(F_j-\Phi(H_j/\sqrt n))
=a_j+c\rho_j\ge a_j+c>\eta$ eventually for some $\eta>0$.
Chebyshev bounds the negative threshold tail by $n/H_j^2$;
the positive discrepancy is at most the normal tail, also at most
$n/H_j^2$. Thus $H_j$ is bounded.

Tightness gives a subsequence $P_j\Rightarrow P$, $H_j\to H$ and
$\rho_j\to b$. The common third-moment ceiling $M$ gives tail
bounds $\int_{|x|>K}x^2dP_j\le M/K$ and
$\int_{|x|>K}|x|dP_j\le M/K^2$. Bounded continuous truncations
therefore preserve mean and variance. Lower semicontinuity gives
$1\le\rho(P)\le b$. For fixed $n$, the sum laws converge weakly.
Applying the closed-set inequality first at $H+t$ and then letting
$t\downarrow0$ gives
$\limsup F_j(H_j)\le F_P(H)$, including an atom at $H$.
Consequently
\[
 a(P,H)\ge A+c\{b-\rho(P)\}\ge A.
\]
Maximality forces equality and $\rho(P)=b$. Convergence of third
moments is a conclusion here; it was not inferred from bounded third
moments. The discrepancy remains positive since $A+c\rho\ge A+c>0$.

For contact, contaminate the raw marginal by
$(1-\epsilon)P+\epsilon\delta_y$. Its mean is $\epsilon y$,
variance $v_\epsilon=1+\epsilon(y^2-1)-\epsilon^2y^2$, and centered
third absolute moment has derivative $|y|^3-\rho-3y\beta$ at zero.
The last assertion follows by dominated differentiation of
$|x-\epsilon y|^3$, with bound $3|y|(|x|+|y|)^2$.
Evaluate the standardized law at the standardized image of the fixed
raw event $S_n\le H$. The raw event probability is a polynomial in
$\epsilon$, with derivative $n(h(y)-F)$ even at an atom collision.
The normal argument is
$(H-n\epsilon y)/\sqrt{nv_\epsilon}$, whose derivative is
$-\sqrt n y-H(y^2-1)/(2\sqrt n)$.
The one-sided derivative of the excess is exactly \eqref{v9:core:contact};
maximality makes it nonpositive.

Its integral against $P$ is zero. All terms are integrable.
The closed-event map $h$ is upper semicontinuous; at an atom it takes
the larger one-sided value. Thus $I$ is upper semicontinuous. Strict
negativity at a closed-support point would hold with a uniform negative
margin on a neighborhood of positive $P$-mass, contradicting its zero
integral. Finally $I(y)\le-c|y|^3+C_2y^2+C_1|y|+C_0$ for constants
depending on this pair. It is strictly negative on both sufficiently
distant rays, proving compact support. No higher moments were used
before this step.
\end{proof}

\subsection{Reflection, equality, and the first unresolved index}
\label{v9:core:reflection}
The sign of the attained maximum detects failure of a strict row bound,
including equality. If $C_n$ denotes the supremum of the normalized
absolute discrepancy, reflection gives, including an atom at $H$,
\[
 \Pp(-S_n\le-H)-\Phi(-z)
 =\Phi(z)-\Pp(S_n<H)\ge\Phi(z)-F.
\]
Thus each negative closed-CDF error is bounded by a positive closed-CDF
error of the reflected law; equality at an atom is not assumed.
Consequently $C_n$ is also the supremum of positive
closed scores $r$. Since $a=\rho(r-c)$, $C_n<c$ gives
$A\le C_n-c<0$, and $C_n>c$ gives $A>0$. If $C_n=c$, choose
$r_j\uparrow c$ with $r_j\ge c/2$. Then
$\rho_j\le2\sqrt n/c$, so $\rho_j(r_j-c)\to0$ and $A=0$.
Hence $C_n\ge c$ supplies an attained pair with $a\ge0$.

\begin{corollary}[First failure of a strict bound]\label{v9:core:first-failure}
Assume a genuine finite base $C_j\le c$ for $1\le j<N$. If a strict
bound $C_n<c$ fails at any index $n\ge N$, there is a least such index,
an attained fixed-target pair with $a\ge0$, and true predecessor bounds
$C_j\le c$ for every $j<n$. The bounds are strict for $N\le j<n$.
\end{corollary}
\begin{proof}
The nonempty set of failing indices at least $N$ has a least member.
The preceding equivalence supplies its maximizing pair. Indices below
$N$ use the finite base; the other earlier indices use minimality.
\end{proof}
Thus an argument excluding every such pair for $n\ge N$ would give
strict bounds above the finite base, including exclusion of equality.
The subsequent entry argument uses only $C_{n-1}\le c$; it does not
use an upper bound for the actual score $c+a/\rho$. The finite base
and the eventual exclusion remain separate obligations.

Equivalently, $\Delta=Q-h\ge0$ everywhere and $\Delta=0$ on the closed
support, where
\begin{equation}\label{v9:core:Q}
 Q(y)=F+\frac{H\varphi+c\rho}{2n^{3/2}}
 -\frac{n\varphi+3c\beta}{n^{3/2}}y
 -\frac{H\varphi+3c\rho}{2n^{3/2}}y^2
 +\frac c{n^{3/2}}|y|^3.
\end{equation}

\subsection{Zero-bias integration and the exact identity}

We now allow an arbitrary compactly supported standardized law, retaining
the definitions of $a,Q,h,\Delta$ above without assuming contact.
Put $\nu=\mathcal L(H-S_{n-1})$ and
$\tau(t)=\E[X\mathbf1_{X\ge t}]$. Then $\tau\ge0$, $\int\tau=1$.
Over the bounded complementary intervals $(u,v)$ of the closed support
$K$, define
\begin{align*}
 J&=\sum\tau_{u,v}(v-u)^2,&
 J_1&=\sum\tau_{u,v}\int_u^v(t-u)t\,dt,\\
 J_2&=\sum\tau_{u,v}\int_u^v(t-u)t|t|\,dt,&
 L&=\sum\tau_{u,v}\int_{[u,v)}(t-u)\,d\nu(t).
\end{align*}
Here $\tau_{u,v}$ is the constant interior value. Introduce two
contact residuals:
\begin{equation}\label{v9:core:residuals}
 \xi=\E[X\Delta(X)],\qquad
 B=\int_K\tau(t)\Delta(t)\,dt
       +\sum_{(u,v)}\tau_{u,v}(v-u)\Delta(v).
\end{equation}
The integral over $K$ is a Lebesgue integral, not an expectation under $P$.
It must be retained when the closed support has positive Lebesgue measure.
If $K\subset[-R,R]$, then $J\le2R$, $|J_1|\le RJ/2$,
$|J_2|\le R^2J/2$, and $0\le L\le2R$. Also $\Delta$ is bounded on
the convex hull, so the sum in $B$ is absolutely convergent because
$\sum\tau_{u,v}(v-u)\le1$.

Let $M_4=\E X^4$, $M_5=\E|X|^5$, and put
\[
 K_q=M_5/4-\rho M_4/2-3\beta\gamma/2+\rho/2,
 \qquad K_1=\E[X|X|^3]-3\beta-3\rho\gamma/2,
\]
\[
 \Omega=J_1+(M_4-3)/6,\quad
 \varepsilon=\E q_z(S_n/\sqrt n)-\E q_z(Z_0)
                  -2\gamma\varphi/(3\sqrt n),\quad
 q_z(x)=(x-z)|x-z|,
\]
where $Z_0$ is standard normal. Define $E_n=-n^{3/2}\varepsilon/2$ and
\begin{equation}\label{v9:core:R}
 \mathcal R_c=c(3\beta J/2+3\rho J_1-3J_2-K_q)
 +H(cK_1+\varphi\Omega)-H^2(c\rho+\gamma\varphi/2)+E_n.
\end{equation}

\begin{theorem}[Conservation with contact residuals]\label{v9:core:general}
For every such law, every $H\in\mathbb R$ and every $c>0$,
\begin{equation}\label{v9:core:offcontact}
 (n+H^2)a+nc\rho-\frac{n\varphi}{6}(\gamma+3J)+n^{3/2}L
 =\mathcal R_c+n^{3/2}(B-H\xi).
\end{equation}
At a fixed-target maximizing pair, $B=\xi=0$. In particular,
\begin{equation}\label{v9:core:conservation}
 (n+H^2)a+nc\rho-\frac{n\varphi}{6}(\gamma+3J)+n^{3/2}L=\mathcal R_c.
\end{equation}
The coefficient of the unnormalized excess is $n+H^2$.
\end{theorem}
\begin{proof}
There are two direct evaluations of $\mathcal A=\int\tau\Delta$.
Since $h(t)=\nu([t,\infty))$, one has
$h(t)=h(v)+\nu([t,v))$ inside a gap. Thus bounded Fubini gives
\[
 \int_u^v\Delta(t)dt=(v-u)\Delta(v)
 +\int_u^v(t-u)(-Q'(t))dt
 -\int_{[u,v)}(t-u)d\nu(t).
\]
An atom at $v$ belongs to $h(v)$ and is not in the last integral;
an atom at $u$ is included but multiplied by zero. Summing, including
the support integral in \eqref{v9:core:residuals}, gives
\begin{equation}\label{v9:core:first}
 \mathcal A=\frac{\varphi J}{2\sqrt n}
 +\frac{3c\beta J/2+(H\varphi+3c\rho)J_1-3cJ_2}{n^{3/2}}-L+B.
\end{equation}

For the second evaluation, signwise Fubini proves
$\int\tau g'=\E[Xg(X)]$ for absolutely continuous $g$ with bounded
derivative on the relevant compact range. In particular
\[
 \int t\tau=\gamma/2,\quad \int t^2\tau=M_4/3,
 \quad\int |t|^3\tau=M_5/4.
\]
Apply that same identity conditionally to
$g_s(x)=\min(s+x,H)$ and then use exchangeability. The choice of
derivative at its single corner is immaterial under the density $\tau$.
It follows directly that
\[
 \int\tau(t)h(t)dt=\E[X_n\min(S_n,H)]
 =\frac1n\E[S_n\min(S_n,H)]
 =T,
\]
where $T=\E[(Z^2-zZ)\mathbf1_{Z\le z}]$, $Z=S_n/\sqrt n$.
The last expression vanishes at an endpoint atom. Integrating
\eqref{v9:core:Q} against $\tau$ therefore yields
\begin{equation}\label{v9:core:second}
 \mathcal A=F-\frac{\gamma\varphi}{2\sqrt n}
 +\frac{H\varphi(3-M_4)/6+cK_q}{n^{3/2}}-T.
\end{equation}

Set $M=\E[Z\mathbf1_{Z\le z}]$. Since
$\E[Xh(X)]=M/\sqrt n$, integration of $XQ(X)$ against $P$ gives
\begin{equation}\label{v9:core:firstmoment}
 M+\varphi=-\frac{\gamma H\varphi}{2n}+\frac{cK_1}{n}-\sqrt n\xi.
\end{equation}
The pointwise identity
$q_z(x)=(x-z)^2-2(x-z)^2\mathbf1_{x\le z}$, with the same identity
for a normal variable, gives
\[
 T-\Phi(z)=-\varepsilon/2-\gamma\varphi/(3\sqrt n)
              +z(M+\varphi)-z^2(F-\Phi(z)).
\]
Inserting \eqref{v9:core:firstmoment} and
$\sqrt n(F-\Phi(z))=c\rho+a$ into \eqref{v9:core:second} gives
\begin{align*}
 \sqrt n\mathcal A={}&c\rho+a-\gamma\varphi/6+P_0/n
                    +\sqrt n\varepsilon/2+H\sqrt n\xi,\\
 P_0={}&H^2(c\rho+a+\gamma\varphi/2)
       -H\{cK_1+\varphi(M_4-3)/6\}+cK_q.
\end{align*}
Equate this with \eqref{v9:core:first} multiplied by $\sqrt n$,
then multiply by $n$ and collect terms. This is
\eqref{v9:core:offcontact}. Under closed-support contact, every
term defining $B$ and $\xi$ vanishes.
\end{proof}

The compensated Fourier representation of this same $E_n$ is proved in
Appendix~\ref{v9:fourier:section}. Its comparator has transform
$e^{-nt^2/2}(1-in\gamma t^3/6)$. The residual identity also explains
why applying the residual-free formula to arbitrary laws would be invalid.
For example, let $X$ be uniform on $[-\sqrt3,\sqrt3]$ and take
$H=0$. There are no support gaps, $F=T=1/2$, $\beta=\gamma=0$,
and $K_q=3\sqrt3/40$. Therefore
$B=\int_K\tau\Delta=cK_q/n^{3/2}>0$.
Omitting the support integral would give a false identity, even though
the gap sums are all empty. This is an off-contact negative control,
not a counterexample to the contact theorem.

\subsection{A positive resource and its ordered phase price}

Put $m=\E X_+=\E X_->0$ and $\sigma=\sqrt{10}$.
Signwise Cauchy gives $a_\pm^2\le mb_\pm$ and
$m^2\le a_+a_-$, where $a_\pm=\E X_\pm^2$ and
$b_\pm=\E X_\pm^3$. Therefore $b_+b_-\ge m^2$, or
$\gamma^2+4m^2\le\rho^2$. A gap on one side of zero has
$(v-u)^2\le2\int_u^v|t|dt$. If a central gap $(-A,B)$ exists,
its extra contribution is $2mAB$, and
$0\le\E[(X+A)(X-B)]=1-AB$. Signwise Tonelli then gives
$J\le\rho+2m$. It follows that
\begin{equation}\label{v9:core:defect}
 D=(3+\sigma)\rho-\gamma-3J
 =3(\rho+2m-J)+(\sigma\rho-\gamma-6m)\ge0.
\end{equation}
The last term is nonnegative by Cauchy in two dimensions.
These statements hold with or without an atom at zero.

Specialize to $c=\CE=(3+\sigma)\phi(0)/6$ and contact. With
$w=H^2/(2n)$, $e=e^{-w}$, the identity becomes
\begin{equation}\label{v9:core:positive}
 \mathcal R_c=(n+H^2)a+nc\rho(1-e)+\frac{n\phi(0)e}{6}D+n^{3/2}L.
\end{equation}
If $a\ge0$, every term is nonnegative. This specialization to $\CE$
is essential for this definition of $D$.

\begin{corollary}[Raw bound, density, and phase price]\label{v9:core:license}\label{v9:core:phase-price}
At such a contact pair with $a\ge0$, assume $n\ge N\ge2$,
$\rho\ge l>0$, $0<c_-\le c$, and a proved raw bound
$\mathcal R_c\le T<Nc_-l$. A negative $T$ excludes the pair.
Otherwise
\begin{gather*}
 \gamma+3J>0,\qquad e\ge1-T/(Nc_-l)>0,\\
 \mathcal U:=(n+H^2)a+n\phi(0)D/6+n^{3/2}L
       \le\mathcal R_c,\\
 \mathcal U\le\mathcal R_c-\tfrac12(c_-l-T/N)H^2.
\end{gather*}
No condition $D<2$ or own-score upper bound is required.
\end{corollary}
\begin{proof}
Set $M_c=c\rho-\phi(0)D/6=\phi(0)(\gamma+3J)/6$.
The general master identity gives
\[
 eM_c=(1+H^2/n)a+c\rho+\sqrt n L-\mathcal R_c/n
        \ge c\rho-T/n>0.
\]
Thus $M_c>0$ and, since $D\ge0$, $M_c\le c\rho$.
The density bound follows. Regrouping \eqref{v9:core:positive} gives
$\mathcal R_c=\mathcal U+n(1-e)M_c$, proving the unweighted bound.
Finally
\[
 n(1-e)M_c=n(e^w-1)eM_c
 \ge\frac{H^2}{2}(c\rho-T/n)
 \ge\frac{H^2}{2}(c_-l-T/N).
\]
This prices one term once. It is a bound on $\mathcal U$, not an
improved bound on the raw remainder.
\end{proof}

Every conclusion holds at each actual index and raw threshold. Common
bounds on a closed moment cell make the result uniform even when the
maximizing law varies with $n$. There is no exchange of a limit with
contact, support gaps, moments, or an integral. In particular, the order
of use is essential: a raw ceiling first proves $M_c>0$ and a density
floor; only then may the Gaussian phase term be subtracted to improve
$\mathcal U$. A smaller resource ceiling is not a replacement for the
raw ceiling in that licensing step.
